\par\bigskip
\SourceLeaf{425}
\begin{center}
\pdfbookmark[2]{Olbers an Gauss, Bremen, 19. December 1810}{pallas-brief-18101219}
\textsc{Olbers} an \textsc{Gauss}, Bremen, 19. December 1810.
\end{center}
\noindent\hspace*{\parindent}\{.\quad.\quad.\quad.\quad.\quad Ich hoffe bald, vielleicht noch diese Woche einen Anlass zu finden, an \textsc{Delambre} zu schreiben und werde, natürlich ohne Sie im geringsten zu compromittiren, zu erfahren suchen, ob nicht der Preistermin für die {\AstroSymbols ⚴}-Perturbationen noch verlängert werden könne.\}
\par\Needspace{7\baselineskip}\bigskip
\begin{center}
\pdfbookmark[2]{Gauss an Schumacher, Göttingen, 6. Januar 1811}{pallas-brief-18110106}
\textsc{Gauss} an \textsc{Schumacher}, Göttingen, 6. Januar 1811.
\end{center}
\noindent\hspace*{\parindent}.\quad.\quad.\quad.\quad.\quad Auf den Fall, dass das Pariser Institut die Pallasstörungen zum viertenmale zur Preisfrage aufgibt, bin ich nicht ganz abgeneigt, darauf zu reflectiren, da jetzt dieser Gegenstand auch an sich Interesse für mich erhalten hat. Ich zweifle aber, dass jenes geschehen wird. Da Sie dort den Moniteur selbst lesen, so bitte ich Sie, wenn Sie in diesen Tagen etwas darüber darin finden, es mir gütigst mitzutheilen.
\par\bigskip
\begin{center}
\pdfbookmark[2]{Olbers an Gauss, Bremen, 26. Januar 1811}{pallas-brief-18110126}
\textsc{Olbers} an \textsc{Gauss}, Bremen, 26. Januar 1811.
\end{center}
\noindent\hspace*{\parindent}\{.\quad.\quad.\quad.\quad.\quad Mit dem grössten Vergnügen habe ich aus dem Moniteur gesehen, dass der Preis für die Perturbationen der Pallas von 6000 Franken bis zum 1. October 1816, doch in dem Masse prorogirt ist, dass er der ersten, während dieses Zeitraums einlaufenden Schrift, die den Absichten und Forderungen der Preisaufgeber entspricht, ertheilt werden soll. Nun zweifle ich nicht\quad.\quad.\quad., dass Sie nach Ihrer mir so erfreulichen Äusserung die vollständige Auflösung dieser Aufgabe unternehmen.\quad.\quad.\quad.\quad.\quad.\}
\par
